\documentclass[french,11pt]{book}
\usepackage{teach}
\usepackage[french]{babel}
\usepackage{amsmath,amssymb}
\usepackage{array}
\newcommand{\R}{\mathbb{R}}
\begin{document}
\begin{center}
{\Large \textbf{Exercices -- Probabilité conditionelle}}\\[2mm]
Terminale technologique
\end{center}

\section*{Exercices}
\begin{enumerate}
  \item Dans une classe, 60\% des élèves pratiquent un sport. Parmi eux, 40\% pratiquent un sport collectif. Calculer la probabilité de choisir un élève qui pratique un sport collectif.
  \item On sait que $P(A)=0,35$ et $P_A(B)=0,6$. Calculer $P(A\cap B)$.
  \item Un test médical est positif chez 98\% des personnes malades et chez 4\% des personnes non malades. La maladie touche 2\% de la population. Construire un arbre pondéré.
  \item Avec l'arbre précédent, calculer la probabilité qu'un test soit positif.
\end{enumerate}

\end{document}
